% 数学ノート集・共通テンプレート v2（積分論2026年度第1部の版面を参考）
% A5判。章内で定義・定理・例・注意等の番号を共通化。
% コンパイル： platex -interaction=nonstopmode math_note_template_v2.tex (2回)
%              dvipdfmx math_note_template_v2.dvi
\documentclass[a5paper,10pt,onecolumn]{jsarticle}
\usepackage[dvipdfmx,top=22mm,bottom=20mm,inner=19mm,outer=18mm,headheight=13pt,headsep=6mm]{geometry}
\usepackage{amsmath,amssymb,amsthm,mathtools}
\usepackage{enumitem,fancyhdr,titlesec}
\usepackage[dvipdfmx,bookmarks=true,colorlinks=false]{hyperref}
\usepackage{pxjahyper}

% ----- 差し替える文書情報 -----
\newcommand{\NoteTitle}{有限体の構造}
\newcommand{\NoteSeries}{数学ノート集}

% ----- 文字・組版 -----
% 参考PDFの質素な明朝＋Computer Modernの印象を引き継ぎながら，
% A5判で窮屈にならないよう本文の行間を少し確保する．
\linespread{1.13}
\setlength{\parindent}{1zw}
\setlength{\parskip}{0pt}
\setlength{\abovedisplayskip}{6pt plus 2pt minus 2pt}
\setlength{\belowdisplayskip}{6pt plus 2pt minus 2pt}
\setlength{\abovedisplayshortskip}{3pt plus 2pt minus 1pt}
\setlength{\belowdisplayshortskip}{4pt plus 2pt minus 1pt}
\setlist[enumerate]{topsep=3pt,itemsep=1pt,parsep=0pt,leftmargin=2.4zw}
\titleformat{\section}{\normalfont\normalsize\bfseries}{\thesection}{0.8zw}{}
\titlespacing*{\section}{0pt}{2.0ex plus .4ex}{.9ex}
\titleformat{\subsection}{\normalfont\normalsize\bfseries}{\thesubsection}{0.7zw}{}
\titlespacing*{\subsection}{0pt}{1.5ex plus .3ex}{.6ex}
\raggedbottom

% ----- 柱とページ番号：ページ上部、罫線なし -----
\fancypagestyle{mathnote}{%
  \fancyhf{}%
  \fancyhead[L]{\footnotesize \NoteSeries\ （\NoteTitle）}%
  \fancyhead[R]{\footnotesize \thepage}%
  \renewcommand{\headrulewidth}{0pt}%
  \renewcommand{\footrulewidth}{0pt}%
}
\pagestyle{mathnote}

% ----- 節内通し番号の定義・定理群 -----
\newtheoremstyle{mathnotetheorem}
 {4.5pt}{4.5pt}{\normalfont}{}%
 {\bfseries}{}{0.65zw}%
 {\thmname{#1}\thmnumber{ #2}\thmnote{（#3）}}
\theoremstyle{mathnotetheorem}
\newtheorem{theorem}{定理}[section]
\newtheorem{definition}[theorem]{定義}
\newtheorem{lemma}[theorem]{補題}
\newtheorem{proposition}[theorem]{命題}
\newtheorem{corollary}[theorem]{系}
\newtheorem{example}[theorem]{例}
\newtheorem{remark}[theorem]{注意}
\renewcommand{\proofname}{証明}
\renewcommand{\qedsymbol}{$\square$}
\makeatletter
\renewenvironment{proof}[1][\proofname]{%
  \par\pushQED{\qed}\normalfont\topsep4\p@\@plus2\p@\relax
  \trivlist\item[\hskip\labelsep\bfseries #1]\ignorespaces
}{\popQED\endtrivlist\@endpefalse}
\makeatother

% ----- 使うノートで適宜拡張する記号 -----
\newcommand{\R}{\mathbb{R}}
\newcommand{\norm}[1]{\lVert #1\rVert}
\newcommand{\abs}[1]{\lvert #1\rvert}

\begin{document}
\section*{読了後にできること}
本稿を読むための前提は，有限体，Frobenius写像，既約多項式．
本稿の具体的な到達課題は，有限体の元を具体的な多項式商環で表し，積・逆元を計算する．

\section{有限体の位数}

本稿では体は可換体とし，有限体の標数・部分体・乗法群を
代数的に特徴づける．
素数 $p$ に対し $\mathbb F_p=\mathbb Z/p\mathbb Z$ と書き，
体 $K$ の零でない元全体の乗法群を $K^\times$ と表す．
体拡大 $L/K$ の次数は $[L:K]=\dim_K L$ とする．

\begin{lemma}\label{lem:char}
有限体 $K$ の標数は素数 $p$ である．
このとき $K$ は素体 $\mathbb F_p$ の有限次拡大であり，
$n=[K:\mathbb F_p]$ とすれば
$|K|=p^n$ である．
\end{lemma}
\begin{proof}
$K$ の標数は有限な正整数 $m$ である．
もし $m=ab$（$1<a,b<m$）なら
$(a1_K)(b1_K)=m1_K=0$ となり，
零因子を持たないという体の性質に反する．
従って $m=p$ は素数である．
$\mathbb Z\to K$ の核は $p\mathbb Z$ なので，
素体 $\mathbb F_p=\mathbb Z/p\mathbb Z$ が $K$ に埋め込まれる．
$K$ は有限集合だから $\mathbb F_p$ 上有限次元である．
その基底を $e_1,\dots,e_n$ とすると，
$K$ の各元は一意に
$\sum_{j=1}^na_je_j$（$a_j\in\mathbb F_p$）
と書けるため，$|K|=p^n$ である．
\end{proof}

\begin{lemma}[Frobenius写像]\label{lem:frob}
標数 $p$ の体 $K$ について
\[
 F:K\to K,\qquad F(x)=x^p
\]
は単射な体準同型である．
$K$ が有限なら $F$ は自己同型である．
\end{lemma}
\begin{proof}
二項定理と
$p\mid{p\choose k}$（$1\le k\le p-1$）
により $(x+y)^p=x^p+y^p$ となる．
また $(xy)^p=x^py^p$，$1^p=1$ なので体準同型である．
体準同型の核は零イデアルなので単射である．
$K$ が有限なら単射は全射でもある．
\end{proof}

\section{存在と一意性}

\begin{theorem}[有限体の存在と一意性]\label{thm:finite-field}
任意の素数 $p$ と正整数 $n$ に対し，
位数 $p^n$ の体が存在する．
また位数 $p^n$ の任意の二つの体は同型である．
\end{theorem}
\begin{proof}
$\mathbb F_p$ 上の多項式
$f(X)=X^{p^n}-X$ の分解体を $L$ とし，
その根全体を
\[
 K=\{x\in L:x^{p^n}=x\}
\]
と置く．
標数 $p$ におけるFrobeniusの加法性から，
$x,y\in K$ なら
$(x-y)^{p^n}=x^{p^n}-y^{p^n}=x-y$ である．
積についても $(xy)^{p^n}=xy$ が成立する．
$x\ne0$ なら
$(x^{-1})^{p^n}=(x^{p^n})^{-1}=x^{-1}$
である．$0,1\in K$ なので $K$ は $L$ の部分体である．
また $f'(X)=p^nX^{p^n-1}-1=-1$ だから
$f$ は重根を持たず，分解体 $L$ 内に相異なる根を
ちょうど $p^n$ 個持つ．従って $|K|=p^n$ である．

一意性について，
$K'$ を位数 $p^n$ の任意の体とする．
乗法群 $(K')^\times$ の位数は $p^n-1$ なので，
Lagrangeの定理から
$x^{p^n-1}=1$（$x\ne0$）である．
従って全ての $x\in K'$ が $x^{p^n}=x$ を満たす．
よって $f$ は $K'$ 上で完全に一次式に分解し，
その根全体は $K'$ に等しい．
従って $K'$ は $f$ の $\mathbb F_p$ 上の分解体である．
分解体の一意性により $K'\cong K$ となる．
\end{proof}

\begin{definition}
定理\ref{thm:finite-field}の一意な同型類を
$\mathbb F_{p^n}$ と書く．
その元の同一視は一般には一意ではなく，
Frobeniusの冪によって異なる自己同型が生じる．
\end{definition}

\section{乗法群の巡回性}

\begin{lemma}\label{lem:order-coprime}
可換群 $G$ の元 $x,y$ の位数がそれぞれ
有限な $m,n$ で $\gcd(m,n)=1$ なら，
$xy$ の位数は $mn$ である．
\end{lemma}
\begin{proof}
可換性により $(xy)^{mn}=1$ である．
$(xy)^k=1$ なら $x^k=y^{-k}$ となる．
左辺は位数が $m$ を割る元，
右辺は位数が $n$ を割る元であるから，
両者の位数は $\gcd(m,n)=1$ を割り，従って双方とも単位元である．
ゆえに $m\mid k$，$n\mid k$ であり，
互いに素であることから $mn\mid k$ となる．
従って位数は $mn$ である．
\end{proof}

\begin{lemma}[有限可換群の最大位数]\label{lem:exponent}
$G$ を有限可換群とし，
$e$ を全ての元の位数の最小公倍数とする．
このとき位数 $e$ の元が存在する．
\end{lemma}
\begin{proof}
$e=\prod_{j=1}^r\ell_j^{a_j}$ を素因数分解する．
各 $j$ について $G$ の元 $g_j$ で
その位数が $\ell_j^{a_j}$ を因数として持つものを取れる．
そうでなければ全ての元の位数の
$\ell_j$-進指数が $a_j$ より小さくなり，
$e$ の定義に矛盾する．
$\operatorname{ord}(g_j)=\ell_j^{a_j}m_j$
かつ $\gcd(\ell_j,m_j)=1$ と書き，
$h_j=g_j^{m_j}$ と置くと
$\operatorname{ord}(h_j)=\ell_j^{a_j}$ である．
$G$ は可換だから，
補題\ref{lem:order-coprime}を繰り返し適用すると
$h=h_1\cdots h_r$ の位数は
$\prod_j\ell_j^{a_j}=e$ となる．
$e=1$ の場合は単位元を取ればよい．
\end{proof}

\begin{theorem}[有限体の乗法群]\label{thm:cyclic}
有限体 $K$ の乗法群 $K^\times$ は巡回群である．
特に $\mathbb F_{p^n}$ には乗法位数 $p^n-1$ の元が存在する．
\end{theorem}
\begin{proof}
$G=K^\times$ と置くと有限可換群である．
補題\ref{lem:exponent}によって位数 $e$ の元が存在する．
$e$ は全ての元の位数の公倍数なので，
全ての $x\in K^\times$ は $x^e=1$ を満たす．
従って $K^\times$ の全ての元は
多項式 $X^e-1$ の根である．体上の次数 $e$ の
非零多項式の根は高々 $e$ 個だから
$|K^\times|\le e$ となる．
一方，位数 $e$ の元を含むから $e\le|K^\times|$ である．
よって $e=|K^\times|$ となり，
位数 $|K^\times|$ の元が $G$ 全体を生成する．
\end{proof}

\begin{corollary}
位数 $q=p^n$ の有限体において，
$X^m=1$ の解の個数は
$\gcd(m,q-1)$ である．
\end{corollary}
\begin{proof}
$K^\times$ の生成元 $g$ を取ると，
任意の元は $g^j$（$0\le j<q-1$）と一意に書ける．
$(g^j)^m=1$ は
$(q-1)\mid mj$ と同値であり，
この合同式の異なる解は
法 $q-1$ で $\gcd(m,q-1)$ 個ある．
\end{proof}

\section{Frobenius自己同型とGalois群}

\begin{theorem}\label{thm:galois}
$K=\mathbb F_{p^n}$ とし，
$F:K\to K$ を $F(x)=x^p$ と定める．
このとき $F$ の位数は $n$ であり，
\[
 \operatorname{Gal}(K/\mathbb F_p)
 =\langle F\rangle\cong\mathbb Z/n\mathbb Z
\]
である．
\end{theorem}
\begin{proof}
$K$ の全ての元は $x^{p^n}=x$ を満たすので
$F^n=\operatorname{id}$ である．
もし $1\le d<n$ について
$F^d=\operatorname{id}$ なら，
$K$ の全ての元が $X^{p^d}-X$ の根となる．
この多項式の次数は $p^d<p^n=|K|$ なので矛盾する．
従って $F$ の位数はちょうど $n$ である．

$K$ は $\mathbb F_p$ 上の
$X^{p^n}-X$ の分解体であり，
導関数が $-1$ なので分離拡大である．
従って $K/\mathbb F_p$ は有限Galois拡大で，
Galois群の位数は
$[K:\mathbb F_p]=n$ に等しい．
$F$ は $\mathbb F_p$ を固定する位数 $n$ の自己同型なので，
Galois群の全体を生成する．
\end{proof}

\begin{theorem}[部分体の分類]\label{thm:subfields}
$K=\mathbb F_{p^n}$ とする．
任意の正の約数 $d\mid n$ に対して，
$K$ は位数 $p^d$ の部分体をただ一つ持つ．
また $K$ の全ての部分体はこの形であり，
$F$ の適切な冪の固定体として記述できる．
\end{theorem}
\begin{proof}
有限Galois拡大 $K/\mathbb F_p$ のGalois群は
$\langle F\rangle$ という位数 $n$ の巡回群である．
$G$ の部分群は約数 $d\mid n$ に対して
$H_d=\langle F^d\rangle$ と一意に対応する．
$H_d$ の位数は $n/d$ なので，
Galois対応による固定体 $K^{H_d}$ の次数は
\[
 [K^{H_d}:\mathbb F_p]
 =\frac{|G|}{|H_d|}=d
\]
となる．従って固定体の位数は $p^d$ である．
Galois対応は包含を逆向きにする全単射なので，
全ての中間体はこのように得られ，
各 $d$ について一意である．
\end{proof}

\begin{corollary}
$\mathbb F_{p^a}$ が $\mathbb F_{p^b}$ に
$\mathbb F_p$-埋め込み可能であることと
$a\mid b$ は同値である．
\end{corollary}
\begin{proof}
埋め込みがあるなら像は $\mathbb F_{p^b}$ の
位数 $p^a$ の部分体なので，定理\ref{thm:subfields}から
$a\mid b$ となる．逆に $a\mid b$ なら定理から
位数 $p^a$ の部分体が存在し，
有限体の一意性によって
$\mathbb F_{p^a}$ と同型である．
\end{proof}

\section{既約多項式とM\"obius反転}

\begin{lemma}\label{lem:irreducible-divisor}
$f\in\mathbb F_p[X]$ を次数 $d\ge1$ のモニック既約多項式とする．
このとき $f\mid X^{p^n}-X$ であることと $d\mid n$ は同値である．
\end{lemma}
\begin{proof}
$f$ の根 $\alpha$ を含む分解体を考える．
$\mathbb F_p(\alpha)$ は位数 $p^d$ の有限体である．
$\alpha^{p^n}=\alpha$ が成り立つことは，
$\alpha$ が $X^{p^n}-X$ の根であることと同値であり，
$f$ の既約性により $f\mid X^{p^n}-X$ と同値である．
一方，$\alpha$ のFrobenius軌道の長さは $d$ である．
実際，$\alpha^{p^d}=\alpha$ なので長さは $d$ を割るが，
長さが $e<d$ なら $\alpha$ は次数 $p^e$ の有限体の元であり，
$\mathbb F_p(\alpha)$ の次数が $d$ であることに矛盾する．
従って $\alpha^{p^n}=\alpha$ であることは
$d\mid n$ と同値となる．
\end{proof}

\begin{theorem}[既約多項式の個数]\label{thm:irreducible-count}
$N_p(n)$ を $\mathbb F_p[X]$ 上の次数 $n$ の
モニック既約多項式の個数とする．
このとき
\[
 p^n=\sum_{d\mid n}d\,N_p(d),
 \qquad
 N_p(n)=\frac1n\sum_{d\mid n}\mu(d)\,p^{n/d}
\]
が成り立つ．
ここで $\mu$ はM\"obius関数である．
\end{theorem}
\begin{proof}
多項式 $X^{p^n}-X$ の導関数は $-1$ なので重根を持たない．
補題\ref{lem:irreducible-divisor}によって，
そのモニック既約因子は次数が $n$ を割るものに限られ，
逆にそのような既約多項式は全て因子となる．
従って
\[
 X^{p^n}-X=
 \prod_{d\mid n}\ \prod_{\substack{f\ {\rm monic\ irreducible}\\\deg f=d}}f(X)
\]
である．両辺の次数を比較すると
$p^n=\sum_{d\mid n}dN_p(d)$ を得る．
この式に整数の約数関係に対する
M\"obius反転公式を適用すれば
$nN_p(n)=\sum_{d\mid n}\mu(d)p^{n/d}$
となる．
\end{proof}

\begin{proposition}[既約多項式の存在]\label{prop:irreducible-exists}
任意の素数 $p$ と正整数 $n$ に対し，
$\mathbb F_p[X]$ には次数 $n$ の既約多項式が存在する．
\end{proposition}
\begin{proof}
$K=\mathbb F_{p^n}$ とする．
定理\ref{thm:cyclic}により
$K^\times$ の生成元 $\alpha$ が存在する．
$\mathbb F_p(\alpha)$ は $\alpha$ を含む部分体であり，
$\alpha$ のべきとして $K^\times$ の全ての元を含むので
$\mathbb F_p(\alpha)=K$ である．
従って $\alpha$ の $\mathbb F_p$ 上の最小多項式は
次数 $[K:\mathbb F_p]=n$ の既約多項式である．
\end{proof}

\section{具体的な有限体の構成}

\begin{example}[$\mathbb F_4$]
$\mathbb F_2[X]$ における $X^2+X+1$ は既約である．
\[
 \mathbb F_4=\mathbb F_2[\alpha],\qquad
 \alpha^2+\alpha+1=0
\]
と置くと，元は $0,1,\alpha,\alpha+1$ の4個であり，
$\alpha^3=1$ なので乗法群は $\alpha$ で生成される．
\end{example}
\begin{proof}
多項式 $X^2+X+1$ は $X=0$，$X=1$ のどちらでも
零にならないので次数2の既約多項式である．
商環 $\mathbb F_2[X]/(X^2+X+1)$ は体であり，
剰余類の各元は $a+b\alpha$（$a,b\in\mathbb F_2$）
と一意に書ける．
$\alpha^2=\alpha+1$ なので
$\alpha^3=\alpha(\alpha+1)=\alpha^2+\alpha=1$ である．
$\alpha\ne1$ だから乗法位数は3となる．
\end{proof}

\begin{example}[$\mathbb F_8$]
$\mathbb F_2[X]$ の多項式 $X^3+X+1$ は既約である．
従って
\[
 \mathbb F_8\cong\mathbb F_2[X]/(X^3+X+1)
\]
であり，$\alpha=X\bmod(X^3+X+1)$ は
$\alpha^3=\alpha+1$ を満たす．
\end{example}
\begin{proof}
次数3の多項式が既約でないなら一次因子を持つ．
$X^3+X+1$ は $0$，$1$ のいずれも根に持たないので
既約である．
商環は次数3の $\mathbb F_2$-拡大だから
元は8個であり，関係式は定義から従う．
\end{proof}

\begin{example}[$\mathbb F_9$]
$\mathbb F_3[X]$ の多項式 $X^2+1$ は既約である．
\[
 \mathbb F_9\cong\mathbb F_3[\beta],
 \qquad \beta^2=-1=2
\]
と置く．
\end{example}
\begin{proof}
$\mathbb F_3$ の平方は $0,1$ なので
$X^2+1$ は根を持たない．
次数2だから既約であり，商は位数9の体である．
\end{proof}

\begin{remark}
有限体は同型を除いて一意だが，その具体的な表示には
既約多項式の選択が必要である．
特に有限体間の埋め込みは，Frobeniusの自己同型のため
一般に一意ではない．
\end{remark}


\section{到達課題の完全解答}
\noindent\textbf{4元体の乗法を実行する．}
\begin{proposition}
$\mathbb F_4=\mathbb F_2[t]/(t^2+t+1)$ は体であり，
$\alpha=[t]$ とすると $\alpha^2=\alpha+1$，
$\alpha^3=1$ が成立する．
\end{proposition}
\begin{proof}
$t^2+t+1$ は $\mathbb F_2$ 上で
$t=0,1$ のいずれでも零とならないので既約である．
従ってその商環は体であり，元は
$0,1,\alpha,\alpha+1$ である．
商の関係式から $\alpha^2=\alpha+1$，
さらに
$\alpha^3=\alpha(\alpha+1)
=\alpha^2+\alpha=1$ を得る．
\end{proof}
\begin{proposition}
$\mathbb F_4^\times$ は $\alpha$ を生成元とする位数3の巡回群である．
\end{proposition}
\begin{proof}
非零元は $1,\alpha,\alpha+1=\alpha^2$ であり，
$\alpha^3=1$，$\alpha\ne1$ だから $\alpha$ の位数は3である．
\end{proof}
\noindent\textbf{到達確認．}
$(\alpha+1)^{-1}$ を求めよ．
\par\noindent\textbf{解答．}
$\alpha+1=\alpha^2$，
$(\alpha^2)\alpha=\alpha^3=1$ だから
逆元は $\alpha$ である．
\end{document}
